\scanpage{53}
\mspage{52}

\uline{13 juin} On n'a pas encore exprimé tous les Axiomes \uline{Cont}. 1, ...., il reste pour compte

\uline{Cont} 3 dim finie

\uline{Cont} 4 ``admissibilité'' -- mais résulte déjà je crois des axiomes C.

\uline{Cont} 6 \inserted{Critère de} disjonction via les lieux -- on y reviendra.

\uline{Cont} 7 Deux lieux de str. dominantes différentes sont disjoints.

Mais avant de continuer, je voudrais encore une fois changer d'optique !

Prendre comme ens de base $\mathcal F$, l'ens des figures, avec les deux relations d'ordre \[ \leq\ \text{(ou l'inclusion)} \qquad \text{et}\qquad \preccurlyeq\ \text{(la subdivision)} \]

\uline{Contrée} (quatrième mouture !)

\unsure{$U$ est|C'est} un ens. $\mathcal F$ d'objets, appelés \uline{figures}, et deux relations d'ordre $\leq$ et $\preccurlyeq$ sur $\mathcal F$ : \[ \mathcal C=(\underline{\mathcal F},\leq,\preccurlyeq) \] satisfaisant à des axiomes\unsure{,|{}} qu'on va énumérer. Si $F$ et $G$ sont deux figures et $F\leq G$, on dit que $F$ est une \uuline{sous}-\uline{figure}, et on écrira aussi souvent $F\subset G$. Si $F\preccurlyeq G$, on dit que $F$ est une \uline{subdivision} de $G$, ou que $F$ subdivise $G$.
